\documentclass[10pt,reqno]{amsart}

\usepackage[T1]{fontenc}
\usepackage[utf8]{inputenc}
\usepackage{newtxtext,newtxmath}
\usepackage{microtype}
\usepackage{mathtools}
\usepackage{enumitem}
\usepackage[nameinlink,noabbrev]{cleveref}
\usepackage[a4paper,margin=28mm,headheight=15pt]{geometry}

\setlength{\parskip}{2pt}
\setlist[enumerate]{leftmargin=7mm,itemsep=2pt,topsep=3pt}
\emergencystretch=3em

\newtheorem{theorem}{Theorem}[section]
\newtheorem{proposition}[theorem]{Proposition}
\newtheorem{lemma}[theorem]{Lemma}
\theoremstyle{definition}
\newtheorem{remark}[theorem]{Remark}

\newcommand{\Q}{\mathbb Q}
\newcommand{\F}{\mathbb F}
\newcommand{\Per}{\operatorname{Per}}

\title[An irrational fixed point in a Salem base]
{An Irrational Fixed Point for a Salem Beta-Transformation:\\
A Scope Correction}

\author{Argus Open Research Project}
\address{Argus Open Research Project}
\date{August 31, 2026}

\subjclass[2020]{Primary 11R06, 37B10; Secondary 11R32}
\keywords{Salem number, beta-transformation, periodic expansion, irrational fixed point, scope correction}

\begin{document}

\begin{abstract}
For the beta-transformation \(T_\beta(x)=\beta x\pmod 1\) on \([0,1)\),
consider
\[
 p(t)=t^4-3t^3+3t^2-3t+1.
\]
We prove that the larger real root \(\beta\) of \(p\) is a Salem number
greater than two and that
\[
 x=\frac1{\beta-1}
\]
is an irrational fixed point of \(T_\beta\).  Hence the literal identity
\(\Per(\beta)=\Q\cap[0,1)\) is false.  This does not address the classical
Salem-base problem formulated by Schmidt, whose right-hand side is
\(\Q(\beta)\cap[0,1)\); indeed the displayed fixed point belongs to
\(\Q(\beta)\).  The result is therefore a correction of an over-restricted
restatement, not a resolution of Schmidt's open problem.
\end{abstract}

\maketitle

\section{Statement and source distinction}

For a real number \(\beta>1\), define
\[
 T_\beta\colon[0,1)\longrightarrow[0,1),
 \qquad T_\beta(x)=\beta x\pmod1.
\]
Following Schmidt \cite{Schmidt}, let \(\Per(\beta)\) be the set of points
whose forward \(T_\beta\)-orbit is finite.  This is the same as eventual
periodicity of the beta-expansion.

A catalogue restatement asserted the following identity for every Salem
number:
\begin{equation}\label{eq:literal}
 \Per(\beta)=\Q\cap[0,1).
\end{equation}
The source does not state \eqref{eq:literal}.  Schmidt proves the corresponding
identity with \(\Q(\beta)\) for Pisot numbers and explicitly leaves open, for
Salem numbers, the question
\begin{equation}\label{eq:schmidt}
 \Per(\beta)=\Q(\beta)\cap[0,1).
\end{equation}

\begin{theorem}\label{thm:main}
Let \(\beta\) be the larger real root of
\[
 p(t)=t^4-3t^3+3t^2-3t+1.
\]
Then \(\beta>2\) is a Salem number, and
\[
 x=\frac1{\beta-1}
\]
belongs to \(\Per(\beta)\setminus\Q\).  Consequently
\eqref{eq:literal} is false.  The point \(x\) belongs to \(\Q(\beta)\), so
the construction is not a counterexample to \eqref{eq:schmidt}.
\end{theorem}

\section{Irreducibility}

\begin{lemma}\label{lem:irreducible}
The polynomial \(p\) is irreducible over \(\Q\).
\end{lemma}

\begin{proof}
Reduction modulo two gives
\[
 \overline p(t)=t^4+t^3+t^2+t+1\in\F_2[t].
\]
It has no root in \(\F_2\), since
\(\overline p(0)=\overline p(1)=1\).  A reducible root-free quartic over
\(\F_2\) would have an irreducible quadratic factor.  The only monic
irreducible quadratic over \(\F_2\) is \(t^2+t+1\).  If \(\alpha\) is one
of its roots, then \(\alpha^3=1\), and
\[
 \overline p(\alpha)
 =\alpha^4+\alpha^3+\alpha^2+\alpha+1
 =\alpha^2\ne0.
\]
Thus \(t^2+t+1\) does not divide \(\overline p\), so \(\overline p\) is
irreducible.  Gauss's lemma proves the claim.
\end{proof}

\section{The Salem conjugates}

For \(t\ne0\), put \(u=t+t^{-1}\).  Then
\begin{equation}\label{eq:reciprocal}
 \frac{p(t)}{t^2}=u^2-3u+1.
\end{equation}
The two values of \(u\) are
\[
 u_+=\frac{3+\sqrt5}{2},
 \qquad
 u_-=\frac{3-\sqrt5}{2},
\]
and therefore
\begin{equation}\label{eq:factorization}
 p(t)=(t^2-u_+t+1)(t^2-u_-t+1)
\end{equation}
over \(\Q(\sqrt5)\).

\begin{proposition}\label{prop:salem}
The four roots of \(p\) are
\[
 \beta,\quad \beta^{-1},\quad z,\quad\overline z,
\]
where \(\beta>2\), \(0<\beta^{-1}<1\), and
\(|z|=|\overline z|=1\).
\end{proposition}

\begin{proof}
Since \(2<\sqrt5<3\),
\[
 u_+>\frac52>2,
 \qquad 0<u_-<\frac12.
\]
The first quadratic in \eqref{eq:factorization} has two positive reciprocal
real roots.  At \(t=2\), its value is
\[
 4-2u_++1=2-\sqrt5<0,
\]
whereas its leading coefficient is positive.  Its larger root \(\beta\)
therefore satisfies \(\beta>2\), and the other is \(\beta^{-1}\).

The discriminant of the second quadratic is negative because
\[
 u_-^2-4<\frac14-4<0.
\]
Its two roots are a nonreal conjugate pair whose product is one; hence each
has modulus one.
\end{proof}

By \Cref{lem:irreducible}, \(p\) is the minimal polynomial of \(\beta\).
It is monic with integer coefficients, so \(\beta\) is an algebraic integer.
\Cref{prop:salem} now proves that \(\beta\) is a Salem number in Schmidt's
sense.

\section{The irrational fixed point}

Set
\[
 x=\frac1{\beta-1}.
\]
Since \(\beta>2\), one has \(0<x<1\).  Moreover,
\begin{equation}\label{eq:fixed}
 \beta x=\frac{\beta}{\beta-1}
 =1+\frac1{\beta-1}=1+x.
\end{equation}
Reducing \eqref{eq:fixed} modulo one gives
\[
 T_\beta(x)=x.
\]
Thus the orbit of \(x\) is the singleton \(\{x\}\), and
\(x\in\Per(\beta)\).

The point is irrational.  Indeed, if \(x\in\Q\), then
\[
 \beta=1+\frac1x\in\Q,
\]
contradicting \Cref{lem:irreducible}.  This proves
\[
 x\in\Per(\beta)\setminus(\Q\cap[0,1)),
\]
and hence proves \Cref{thm:main}.  On the other hand,
\(x=1/(\beta-1)\in\Q(\beta)\), so the same point is compatible with
Schmidt's open identity \eqref{eq:schmidt}.

\section{Exact computational check}

The proof above is entirely symbolic.  An accompanying exact-arithmetic
check independently performs the following bounded tasks:

\begin{enumerate}
\item reduce \(p\) modulo two and exhaust the possible linear and irreducible
  quadratic factors;
\item verify the reciprocal identity \eqref{eq:reciprocal} and reconstruct
  \eqref{eq:factorization};
\item check the strict rational bounds used to obtain \(\beta>2\) and the
  unit-circle conjugate pair; and
\item verify the polynomial identity obtained by eliminating \(\beta\) from
  \(x=1/(\beta-1)\), together with the fixed-point relation
  \eqref{eq:fixed}.
\end{enumerate}

A separate typed polynomial-factor operation returns \(p\) as one monic
factor and verifies exact product reconstruction.  Since that operation does
not itself certify factor irreducibility, the mathematical conclusion uses
the explicit modulo-two proof in \Cref{lem:irreducible}.

\section{Scope and referee checklist}

The result corrects only the right-hand side \(\Q\) in the literal catalogue
identity.  It does not resolve or refute Schmidt's \(\Q(\beta)\) problem.
The distinction is substantive, not cosmetic: the counterexample point is
itself an element of \(\Q(\beta)\).

An external referee can check the conclusion by the following short list.

\begin{enumerate}
\item Verify irreducibility of \(\overline p\) over \(\F_2\).
\item Verify the substitution \(u=t+t^{-1}\) and the two quadratic factors.
\item Check the real-root and unit-circle inequalities without numerical
  approximation.
\item Check \(0<x<1\), \(\beta x=1+x\), and the irrationality argument.
\item Compare the two distinct source formulations
  \(\Q\cap[0,1)\) and \(\Q(\beta)\cap[0,1)\).
\end{enumerate}

\begin{thebibliography}{9}

\bibitem{Schmidt}
K. Schmidt,
On periodic expansions of Pisot numbers and Salem numbers,
\emph{Bull. London Math. Soc.} \textbf{12} (1980), no.~4, 269--278.

\end{thebibliography}

\end{document}
